% Cette trace écrite à été développée par
% + Héloïse Troublé <heloise.trouble@ens-rennes.fr>
% + Erwan Tanguy-Legac <erwan.tanguy-legac@ens-rennes.fr>

% Elle est distribuée sous la licence Creative-Commons by-nc-sa 4.0
% (Attribution, Non-Commercial, Share Alike)
% qui peut être trouvée dans
% [le site de Creative Commons](https://creativecommons.org/licenses/by-nc-sa/4.0/)

\documentclass[11pt,a5paper]{article}
\usepackage[left=1.5cm, right=1.5cm, lines=45, top=0.3in, includehead]{geometry}
\usepackage[T1]{fontenc}
\usepackage[french]{babel}
\usepackage[utf8]{inputenc}
\usepackage{fancyhdr}
\usepackage{titling}
\usepackage{ragged2e}
\usepackage{subcaption}
\usepackage{wrapfig}
\usepackage{tikz}

\DeclareCaptionFormat{custom}{#3}
\captionsetup{format=custom}

\setlength{\droptitle}{-2em}

\predate{}
\postdate{}
\date{\vspace{-12ex}}

\lhead{Informatique sans ordi}
\rhead{29 mars 2024}
\title{\vspace{-5ex}Marmottes}

\pagenumbering{gobble}

\begin{document}

\maketitle
\thispagestyle{fancy}
\justifying

C’est l’hiver, les marmottes s’apprêtent à hiberner. Elles veulent pour cela construire un terrier. Afin qu'il ne s’effondre pas, elles cons-truisent des galeries en forme de V inversé uniquement. Les marmottes ne peuvent pas dormir au
milieu d'une galerie, et dès qu'une marmotte remonte à la surface, elle fait du bruit pour chaque galerie qu'elle traverse. On veut construire une galerie afin que le moins bruit possible soit généré.
Dans l'exemple ci-dessous, les marmottes sont représentées par leur nombre de réveils. 
\section*{Construction de la galerie}

\begin{wrapfigure}{r}{0.5\textwidth}
\begin{tikzpicture}[every node/.style={circle, draw}]
  \node (1) {1};
  \node (2) [right of=1] {2};
  \node (2bis) [right of=2] {2};
  \node (3) [right of=2bis]{3};
  \node (4) [right of=3]{4};
\end{tikzpicture}
\caption{Situation de départ : cercles triés}
\end{wrapfigure}

L'algorithme peut être décrit ainsi : tant qu'il reste au moins deux galeries en construction, on relie les deux plus petites.
La taille d'une galerie est définie comme la somme des chiffres présents dans ses cercles.

\begin{figure}[h!]
\begin{subfigure}{0.5\textwidth}
  \centering
\begin{tikzpicture}[every node/.style={circle, draw}]
  \node at (0, 0) (2)  {2};
  \node at (1, 0) (1) {1};
  \node at (2, 0) (2bis) {2};
  \node at (3, 0) (3) {3};
  \node at (4, 0) (4) {4};
  \node at (1.5, 0.7) (5) {};
  \draw (1) -- (5);
  \draw (2bis) -- (5);
\end{tikzpicture}
  \caption{On relie les cercles à gauche et on trie}
\end{subfigure}%
\hfill
\begin{subfigure}{0.5\textwidth}
  \centering
\begin{tikzpicture}[every node/.style={circle, draw}]
  \node at (3, 0.7) (2)  {2};
  \node at (3.5, 0) (1) {1};
  \node at (4.5, 0) (2bis) {2};
  \node at (1, 0) (3) {3};
  \node at (2, 0) (4) {4};
  \node at (4, 0.7) (5) {};
  \node at (3.5, 1.4) (6) {};
  \draw (1) -- (5);
  \draw (2bis) -- (5);
  \draw (2) -- (6);
  \draw (5) -- (6);
\end{tikzpicture}
  \caption{\hspace{3mm}On continue : la galerie droite vaut 5}
\end{subfigure}
\end{figure}

\begin{figure}[h!]
\begin{subfigure}{0.5\textwidth}
  \centering
\begin{tikzpicture}[every node/.style={circle, draw}]
  \node at (3, 0.7) (2)  {2};
  \node at (3.5, 0) (1) {1};
  \node at (4.5, 0) (2bis) {2};
  \node at (5.5, 0) (3) {3};
  \node at (6.5, 0) (4) {4};
  \node at (4, 0.7) (5) {};
  \node at (3.5, 1.4) (6) {};
  \node at (6, 0.7) (7) {};
  \draw (1) -- (5);
  \draw (2bis) -- (5);
  \draw (2) -- (6);
  \draw (5) -- (6);
  \draw (3) -- (7);
  \draw (4) -- (7);
\end{tikzpicture}
  \caption{On relie les deux galeries à gauche et on trie. La galerie obtenue est de taille 7}
\end{subfigure}%
\hfill
\begin{subfigure}{0.5\textwidth}
  \centering
\begin{tikzpicture}[every node/.style={circle, draw}]
  \node at (3, 0.7) (2)  {2};
  \node at (3.5, 0) (1) {1};
  \node at (4.5, 0) (2bis) {2};
  \node at (5, 0.7) (3) {3};
  \node at (6, 0.7) (4) {4};
  \node at (4, 0.7) (5) {};
  \node at (3.5, 1.4) (6) {};
  \node at (5.5, 1.4) (7) {};
  \node at (4.5, 2.1) (root) {};
  \draw (1) -- (5);
  \draw (2bis) -- (5);
  \draw (2) -- (6);
  \draw (5) -- (6);
  \draw (3) -- (7);
  \draw (4) -- (7);
  \draw (7) -- (root);
  \draw (6) -- (root);
\end{tikzpicture}
  \caption{Galerie finale obtenue}
\end{subfigure}
\end{figure}

Cet algorithme est utilisé pour coder des informations afin qu'elles prennent le moins de place possible dans l'ordinateur.

\end{document}
